\documentclass[11pt,a4paper]{article}

% ---------------------------------------------------------
% PREAMBLE (stable, warning-free, Overleaf-safe)
% ---------------------------------------------------------
\usepackage[utf8]{inputenc}
\usepackage[T1]{fontenc}
\usepackage{lmodern}
\usepackage{amsmath, amssymb, amsfonts}
\usepackage{bm}
\usepackage{geometry}
\usepackage{graphicx}
\usepackage{hyperref}
\usepackage{cite}
\usepackage{authblk}
\usepackage{physics}
\usepackage{mathtools}

\geometry{margin=2.4cm}

\title{\textbf{Companion LIX: Mass-Mapping Inversion in the Relaxation-Driven Cyclic Cosmology}}
\author{Michael Lehmann \\ \texttt{mi.lehmann@gmx.de}}
\date{April 2026}

\begin{document}
\maketitle

\begin{abstract}
Mass-mapping techniques reconstruct the projected mass density or gravitational 
potential from weak-lensing shear fields. In the standard cosmological model, the 
inversion from shear to convergence and potential is determined by the lensing 
kernel, the geometry of the Universe, and the statistical properties of the shear 
field. In the Relaxation-Driven Cyclic Cosmology (RDCC), the scale-dependent 
evolution of the gravitational potential modifies the inversion kernel, the 
reconstruction response, and the coherence of the reconstructed mass maps in a 
characteristic way. This Companion develops a continuous description of mass-
mapping inversion in RDCC, deriving how the relaxation scale influences the 
reconstruction of $\kappa$ and $\Phi$, the mode coupling, and the observational 
signatures in lensing surveys. The resulting framework provides a distinctive probe 
of the RDCC gravitational field in real-space mass maps.
\end{abstract}
% ---------------------------------------------------------
\section{Introduction}

Weak gravitational lensing distorts the shapes of background galaxies, allowing the 
reconstruction of the projected mass density (convergence $\kappa$) or the 
gravitational potential $\Phi$. Mass-mapping techniques invert the shear field to 
obtain $\kappa$ or $\Phi$, providing a powerful probe of the matter distribution 
and the gravitational field.

In the standard cosmological model, the inversion is determined by the lensing 
kernel, the geometry of the Universe, and the statistical properties of the shear 
field. In RDCC, the gravitational potential evolves differently. The relaxation 
mechanism suppresses the decay of small-scale modes and enhances the coherence of 
large-scale modes. This modifies the inversion kernel, the reconstruction response, 
and the coherence of the reconstructed mass maps in a scale-dependent manner, 
producing a distinctive signature in mass-mapping analyses.
% ---------------------------------------------------------
\section{The RDCC Convergence Field}

The convergence is related to the gravitational potential by
\begin{equation}
    \kappa(\hat{\mathbf{n}})
    = \frac{1}{2}\nabla^{2}\phi(\hat{\mathbf{n}}).
\end{equation}

In RDCC, the lensing potential evolves as
\begin{equation}
    \phi_{\mathrm{RDCC}}
    = \phi
      \left[
        1 - \chi_{\phi}
        \left(\frac{\ell}{\ell_{\mathrm{rel}}}\right)^{\alpha}
      \right].
\end{equation}

The RDCC convergence becomes
\begin{equation}
    \kappa_{\mathrm{RDCC}}(\ell)
    = \kappa(\ell)
      \left[
        1 - \chi_{\kappa}
        \left(\frac{\ell}{\ell_{\mathrm{rel}}}\right)^{\alpha}
      \right].
\end{equation}

This modification reflects the scale-dependent evolution of the gravitational 
potential.
% ---------------------------------------------------------
\section{Mass-Mapping Inversion in RDCC}

The standard inversion from shear to convergence is
\begin{equation}
    \hat{\kappa}(\ell)
    = D(\ell)\, \gamma(\ell),
\end{equation}
where $D(\ell)$ is the Kaiser–Squires kernel.

In RDCC, the shear field becomes
\begin{equation}
    \gamma_{\mathrm{RDCC}}(\ell)
    = \gamma(\ell)
      \left[
        1 - \chi_{\gamma}
        \left(\frac{\ell}{\ell_{\mathrm{rel}}}\right)^{\alpha}
      \right].
\end{equation}

The reconstructed convergence becomes
\begin{equation}
    \hat{\kappa}_{\mathrm{RDCC}}(\ell)
    = D(\ell)\,
      \gamma(\ell)
      \left[
        1 - \chi_{\gamma}
        \left(\frac{\ell}{\ell_{\mathrm{rel}}}\right)^{\alpha}
      \right].
\end{equation}

This modification alters the reconstruction response.
% ---------------------------------------------------------
\section{Reconstruction Response and RDCC Signatures}

The reconstruction response is defined as
\begin{equation}
    R_{\ell}
    = \frac{\langle \hat{\kappa}_{\mathrm{RDCC}}(\ell) \rangle}
           {\kappa_{\mathrm{RDCC}}(\ell)}.
\end{equation}

In RDCC, this becomes
\begin{equation}
    R_{\ell,\mathrm{RDCC}}
    = R_{\ell}
      \left[
        1 - \chi_{R}
        \left(\frac{\ell}{\ell_{\mathrm{rel}}}\right)^{\alpha}
      \right].
\end{equation}

The relaxation mechanism enhances large-scale coherence and suppresses small-scale 
structure. Mass maps therefore exhibit:

\begin{itemize}
    \item enhanced large-scale potential wells,
    \item suppressed small-scale peaks,
    \item a characteristic turnover near $\ell_{\mathrm{rel}}$,
    \item modified reconstruction noise at high $\ell$.
\end{itemize}
% ---------------------------------------------------------
\section{Real-Space Mass Maps in RDCC}

The real-space convergence map is
\begin{equation}
    \kappa(\hat{\mathbf{n}})
    = \int \frac{d^{2}\ell}{(2\pi)^{2}}\,
      \kappa(\ell)\,
      e^{i\ell\cdot\hat{\mathbf{n}}}.
\end{equation}

In RDCC, the reconstructed map becomes
\begin{equation}
    \kappa_{\mathrm{RDCC}}(\hat{\mathbf{n}})
    = \kappa(\hat{\mathbf{n}})
      - \chi_{\kappa}
        \int \frac{d^{2}\ell}{(2\pi)^{2}}\,
        \kappa(\ell)
        \left(\frac{\ell}{\ell_{\mathrm{rel}}}\right)^{\alpha}
        e^{i\ell\cdot\hat{\mathbf{n}}}.
\end{equation}

This modification produces a distinctive smoothing pattern tied to the relaxation 
scale.
% ---------------------------------------------------------
\section{Conclusion}

Mass-mapping inversion in the Relaxation-Driven Cyclic Cosmology reflects the 
scale-dependent evolution of the gravitational potential and the matter 
distribution. The relaxation mechanism modifies the inversion kernel, the 
reconstruction response, and the coherence of the reconstructed mass maps in a 
scale-dependent manner, producing distinctive signatures in weak-lensing surveys. 
These effects provide a direct observational window into the relaxation scale and 
the temporal structure of the RDCC gravitational field. The present Companion 
establishes the theoretical foundation for future studies of mass reconstruction, 
nonlinear structure, and the interplay between light and gravity in RDCC.

\bibliographystyle{apsrev4-2}
\nocite{*}
\bibliography{references}

\end{document}
